\clearpage
\begin{examplebox}
\textbf{\textcolor{CExample}{例 1（由坐标求距离）}}
在 $y^2=4x$ 上，$P$ 的横坐标为 $3$，求 $PF$。
\textbf{解：}$2p=4$，故 $p/2=1$，$PF=3+1=\boxed{4}$。
\textbf{回验：}$y_0^2=12$，用两点距离得 $\sqrt{(3-1)^2+12}=4$。

\medskip
\textbf{\textcolor{CExample}{例 2（由距离反求坐标）}}
在 $y^2=8x$ 上，$PF=6$，求 $x_0$。
\textbf{解：}$p/2=2$，$x_0+2=6$，所以 $x_0=\boxed{4}$。
\textbf{回验：}$y_0^2=32$，$\sqrt{(4-2)^2+32}=6$，且 $x_0\ge0$。

\medskip
\textbf{\textcolor{CExample}{例 3（焦半径转坐标，再求面积）}}
在 $y^2=4x$ 上，$PF=5$，求 $\triangle OFP$ 的面积。
\begin{center}% Exact example 3: P=(4,4), F=(1,0), triangle area 2.
\begin{tikzpicture}[>=stealth,font=\small]
\draw[rounded corners=3pt,gray!45] (0,0) rectangle (14.2,8.7);
\node at (7.1,8.15) {例 3：先由 $PF=5$ 求坐标，再找 $\triangle OFP$ 的底和高};
\begin{scope}[shift={(2.2,2.3)},scale=1.05]
 \fill[blue!12] (0,0)--(1,0)--(4,4)--cycle;
 \draw[->,gray] (-.6,0)--(5.2,0) node[right] {$x$};
 \draw[->,gray] (0,-.5)--(0,4.8) node[above] {$y$};
 \draw[teal!75!black,thick,domain=0:4.25,samples=90] plot ({\x*\x/4},\x);
 \draw[blue!75!black,thick] (0,0)--(4,4)--(1,0);
 \draw[blue!75!black,very thick] (0,0)--(1,0);
 \draw[violet,thick,dashed] (4,0)--(4,4) node[midway,right=4pt] {高 $4$};
 \draw[violet] (3.75,0)--(3.75,.25)--(4,.25);
 \fill (0,0) circle(2pt) node[below left=3pt] {$O$};
 \fill (1,0) circle(2pt) node[below=4pt] {$F(1,0)$};
 \fill (4,0) circle(2pt) node[below=4pt] {$H(4,0)$};
 \fill (4,4) circle(2pt) node[right=4pt] {$P(4,4)$};
 \draw[<->,blue!75!black] (0,-.9)--(1,-.9) node[midway,below] {底 $1$};
\end{scope}
\node[anchor=west,align=left] at (9.4,5.1) {$x_0=5-1=4$\\[8pt]$|y_0|=4$\\[10pt]$\displaystyle S=\frac12\cdot1\cdot4=2$};
\node at (7.1,.55) {高落在底边的延长线上也成立；取 $P(4,-4)$，面积仍为 $2$。};
\end{tikzpicture}
\end{center}
\textbf{解：}$F(1,0)$；由 $x_0+1=5$ 得 $x_0=4$。
代回 $y_0^2=4x_0=16$，得 $y_0=\pm4$。
底 $OF=1$，高 $|y_0|=4$，所以 $S=\frac12\cdot1\cdot4=\boxed{2}$。
\textbf{回验：}$PF=\sqrt{(4-1)^2+4^2}=5$。
\end{examplebox}
\clearpage
\begin{examplebox}
  \medskip
  \textbf{\textcolor{CExample}{例 4（换开口，直接识别负号）}}

  \textbf{题目：}已知点 $P(x_0,y_0)$ 在抛物线 $y^2=-8x$ 上，且 $x_0=-3$，求 $PF$。

  \textbf{最短解法：}向左开口，$2p=8$，所以 $p/2=2$，
  \[
    PF=-x_0+\frac p2=-(-3)+2=\boxed{5}.
  \]
  \textbf{校验：}焦点 $F(-2,0)$，$y_0^2=-8(-3)=24$，
  故 $PF=\sqrt{(-3+2)^2+24}=\sqrt{25}=5$。

  \medskip
  \textbf{\textcolor{CExample}{例 5（上下对称，一题两式）}}

  \textbf{题目：}在 $x^2=6y$ 上点 $P$ 的纵坐标为 $5$；
  在 $x^2=-6y$ 上点 $Q$ 的纵坐标为 $-5$。
  分别求两点到各自抛物线焦点的距离。

  \textbf{最短解法：}$2p=6$，$p/2=3/2$。
  向上用 $y_0$，向下用 $-y_0$：
  \[
    PF=5+\frac32=\boxed{\frac{13}{2}},\qquad
    QF'=-(-5)+\frac32=\boxed{\frac{13}{2}}.
  \]
  \textbf{理解：}上下只是方向相反，点沿开口方向离顶点都走了 $5$。

  \medskip
  \textbf{\textcolor{CExample}{例 6（考场极值，一行出答案）}}

  \textbf{题目：}抛物线 $y^2=12x$ 上一点 $P$ 的焦半径 $PF$ 最小是多少？

  \textbf{解法：}$p=6$，$x_0\ge0$，所以
  \[
    PF=x_0+3\ge\boxed{3}.
  \]
  当 $P=O(0,0)$ 时取等号。\textbf{检查：}顶点到焦点的距离的确为 $3$。

  \medskip
  \textbf{\textcolor{CExample}{例 7（逆向思考：先核验点在不在曲线上）}}

  \textbf{题目：}对 $y^2=4x$，若随意取平面内点 $R(3,0)$，能否说 $RF=3+1=4$？

  \textbf{解法：}不能。$R$ 不在曲线上（$0^2\ne4\cdot3$）。
  焦点 $F(1,0)$，真正的距离是 $RF=|3-1|=\boxed{2}$。
  \textbf{结论：}焦半径坐标公式\emph{必须先满足“点在抛物线上”}，不能把它当作任意点的距离公式。
\end{examplebox}
